\chapter{Nền tảng lý thuyết}
\label{chap:c2_theory}

\section{Kiến trúc Modular Monolith cho thương mại điện tử}
\label{sec:c2_ecommerce_overview}

Thiết kế hệ thống thương mại điện tử hoàn chỉnh đòi hỏi phối hợp nhiều thành phần: danh mục sản phẩm, giỏ hàng đồng bộ giữa guest và user, đặt hàng với tính toàn vẹn ACID, và phân quyền đa tác nhân~\cite{turban2018electronic}. \textbf{Modular Monolith}~\cite{richter2020modular} là kiến trúc cân bằng giữa monolithic và microservices: triển khai như một ứng dụng đơn nhất nhưng tổ chức nội bộ thành các module với ranh giới rõ ràng, giao tiếp qua EventBus, shared models và shared services. Kiến trúc này phù hợp với giai đoạn phát triển ban đầu và có thể chuyển sang microservices khi cần.

\section{Các phương pháp chatbot và khoảng trống kỹ thuật}
\label{sec:c2_chatbot_survey}

\textbf{Chatbot rule-based}~\cite{adamopoulou2020chatbots} phản hồi dựa trên tập luật cố định, phù hợp cho tác vụ cấu trúc (tra cứu đơn hàng, chính sách đổi trả) nhưng không hiểu ngôn ngữ tự nhiên và không truy vấn được catalog thực tế~\cite{folstad2017chatbots}. \textbf{Chatbot ML truyền thống}~\cite{chen2021conversational} cải thiện qua intent classification và NER, nhưng vẫn trả lời từ kho mẫu cố định, không tổng hợp được thông tin đa nguồn và cần hàng nghìn mẫu gán nhãn. Bảng~\ref{tab:chatbot_comparison} tổng hợp so sánh.

\begin{table}[H]
  \centering
  \caption{So sánh các phương pháp chatbot theo tiêu chí tư vấn sản phẩm}
  \label{tab:chatbot_comparison}
  \begin{tabular}{|p{4cm}|p{2.5cm}|p{2.5cm}|p{3.5cm}|}
    \hline
    \textbf{Tiêu chí} & \textbf{Rule-based} & \textbf{ML truyền thống} & \textbf{RAG + LLM (TechStore)} \\
    \hline
    Hiểu câu hỏi tự nhiên & Không & Hạn chế & Có \\
    \hline
    Truy xuất dữ liệu catalog thực tế & Không & Không & Có \\
    \hline
    Sinh câu trả lời linh hoạt & Không & Không & Có \\
    \hline
    Duy trì ngữ cảnh đa lượt & Không & Hạn chế & Có \\
    \hline
    Bảo trì khi catalog thay đổi & Thủ công & Thủ công & Tự động \\
    \hline
  \end{tabular}
\end{table}

Cả hai thế hệ chatbot truyền thống đều không đáp ứng yêu cầu cốt lõi: trả lời câu hỏi kỹ thuật tự do dựa trên catalog thực tế. RAG kết hợp LLM lấp đầy khoảng trống này.

\section{Kỹ thuật RAG (Retrieval-Augmented Generation)}
\label{sec:c2_rag}

Phần này trình bày kỹ thuật RAG từ nền tảng LLM, nguyên lý hoạt động, đến Hybrid Search và các thách thức khi áp dụng cho tiếng Việt.

\subsection{Mô hình ngôn ngữ lớn và giới hạn kiến thức}

Mô hình ngôn ngữ lớn (Large Language Model -- LLM) là các mô hình học sâu được huấn luyện trên tập dữ liệu văn bản khổng lồ với hàng tỷ tham số, có khả năng hiểu và sinh văn bản tự nhiên trong nhiều ngữ cảnh phức tạp. Các mô hình tiêu biểu bao gồm GPT-4 (OpenAI), Claude (Anthropic), Gemini (Google) và các mô hình mã nguồn mở như LLaMA và Mistral~\cite{brown2020language}.

Tuy nhiên, LLM có một hạn chế cơ bản: kiến thức của mô hình bị đóng băng tại thời điểm huấn luyện (knowledge cutoff) và không thể cập nhật theo thời gian thực. Đây là vấn đề nghiêm trọng trong các ứng dụng thực tế như tư vấn sản phẩm, khi giá cả, tình trạng tồn kho và danh mục sản phẩm thay đổi liên tục. Ngoài ra, LLM còn có xu hướng ``ảo giác'' (hallucination), tức hiện tượng sinh ra thông tin nghe có vẻ hợp lý nhưng thực tế sai, khi không có dữ liệu tham chiếu đủ cụ thể~\cite{ji2023survey}.

\subsection{Nguyên lý RAG}

RAG (Retrieval-Augmented Generation) là kỹ thuật được giới thiệu bởi Lewis và cộng sự nhằm giải quyết các hạn chế của LLM thuần túy~\cite{lewis2020rag}. Trước khi pipeline hoạt động, dữ liệu nguồn (sản phẩm, chính sách, FAQ) phải được xử lý và vector hóa offline để lưu vào vector store; bước chuẩn bị này thường được gọi là \textit{Indexing}. Tại runtime, pipeline gồm ba giai đoạn chính.

\textbf{Giai đoạn Retrieval} (truy xuất): Khi nhận câu hỏi từ người dùng, query được vector hóa và so sánh với các vector trong store để tìm ra các tài liệu liên quan nhất theo độ tương đồng ngữ nghĩa.

\textbf{Giai đoạn Augmentation} (tăng cường ngữ cảnh): Các tài liệu truy xuất được ghép vào câu hỏi gốc để xây dựng prompt đầy đủ gửi lên LLM. Đây là bước tạo ra sự khác biệt cốt lõi của RAG: thay vì chỉ dựa vào tham số đã học, LLM nhận được ngữ cảnh thực tế cụ thể trong mỗi lần suy luận.

\textbf{Giai đoạn Generation} (sinh văn bản): LLM tiếp nhận prompt đã được tăng cường và sinh phản hồi dựa trên dữ liệu tham chiếu được cung cấp, giảm thiểu hiện tượng ảo giác (hallucination) so với LLM không có retrieval.

So với fine-tuning (huấn luyện lại mô hình trên dữ liệu miền), RAG có ưu điểm không đòi hỏi tài nguyên tính toán lớn, dễ cập nhật khi dữ liệu thay đổi, có thể trích dẫn nguồn gốc thông tin, và dễ kiểm soát chất lượng~\cite{gao2023retrieval}.

Các nghiên cứu gần đây phân loại hệ thống RAG thành ba thế hệ~\cite{gao2023retrieval}. \textit{Naive RAG} là pipeline cơ bản gồm đúng ba bước index, retrieve và generate mà không có tối ưu thêm, dễ triển khai nhưng dễ bị ảnh hưởng bởi chất lượng query đầu vào. \textit{Advanced RAG} bổ sung các bước tiền xử lý trước retrieval (chuẩn hóa query, phân loại ý định, kiểm tra bảo mật) và hậu xử lý sau retrieval (hợp nhất kết quả, reranking, cơ chế dự phòng). \textit{Modular RAG} đẩy xa hơn bằng kiến trúc plug-and-play nơi từng thành phần có thể thay thế độc lập. Pipeline RAG của TechStore thuộc nhóm \textit{Advanced RAG}: trước retrieval có bước chuẩn hóa ngôn ngữ, phân loại ý định và kiểm tra bảo mật bằng regex; sau retrieval có hợp nhất Hybrid Search với overlap boost, đánh dấu \texttt{lowConfidence} cho kết quả kém tin cậy và cơ chế dự phòng từ khóa khi không có kết quả đạt ngưỡng cosine similarity.

\subsection{Vector Embedding và không gian ngữ nghĩa}

\textbf{Vector embedding} là phương pháp ánh xạ văn bản sang không gian vector số thực nhiều chiều, sao cho các văn bản có nội dung tương đồng được biểu diễn bằng các vector gần nhau theo độ đo khoảng cách cosine~\cite{mikolov2013word2vec}. Chất lượng của embedding quyết định trực tiếp chất lượng của bước retrieval.

Sentence-BERT~\cite{reimers2019sentencebert} là mô hình embedding tiêu biểu cho văn bản đơn ngữ, với số chiều phụ thuộc vào kiến trúc cụ thể của từng biến thể mô hình. Mô hình multilingual-e5-large của HuggingFace~\cite{wang2022text} mở rộng khả năng này sang đa ngôn ngữ trong không gian vector 1024 chiều, phù hợp với dữ liệu sản phẩm tiếng Việt lẫn tiếng Anh. Mô hình Jina Embeddings v3~\cite{sturua2024jina} cũng xuất vector 1024 chiều, được tối ưu hóa đặc biệt cho tác vụ retrieval với cơ chế attention cải tiến.

Độ tương đồng giữa hai vector thường được đo bằng \textbf{cosine similarity}:

\begin{equation}
\text{cosine\_sim}(\mathbf{u}, \mathbf{v}) = \frac{\mathbf{u} \cdot \mathbf{v}}{\|\mathbf{u}\| \|\mathbf{v}\|}
\label{eq:cosine_similarity}
\end{equation}

Giá trị cosine similarity nằm trong khoảng $[-1, 1]$, với giá trị gần 1 cho thấy hai văn bản tương đồng cao về ngữ nghĩa. Trong TechStore, ngưỡng mặc định là 0.45; các kết quả có cosine similarity dưới ngưỡng này bị loại trừ khỏi ngữ cảnh gửi cho LLM.

\subsection{Hybrid Search -- kết hợp tìm kiếm ngữ nghĩa và từ khóa}
\label{sec:c2_hybrid_search}

Tìm kiếm ngữ nghĩa thuần túy (dense retrieval) đôi khi bỏ sót các tài liệu khớp chính xác về mặt từ ngữ nhưng có vector không thực sự gần nhau. Ngược lại, tìm kiếm từ khóa truyền thống (sparse retrieval) không nắm bắt được ngữ nghĩa. Để khắc phục, \textbf{Hybrid Search} kết hợp cả hai phương pháp~\cite{ma2021replication}.

BM25 (Best Matching 25)~\cite{robertson1994okapi} là thuật toán tìm kiếm từ khóa hiệu quả, tính điểm dựa trên tần suất từ (TF) và tần suất nghịch đảo tài liệu (IDF) với chuẩn hóa độ dài:

\begin{equation}
\text{BM25}(q, d) = \sum_{t \in q} \text{IDF}(t) \cdot \frac{tf(t,d) \cdot (k_1 + 1)}{tf(t,d) + k_1 \cdot \left(1 - b + b \cdot \frac{|d|}{\overline{dl}}\right)}
\label{eq:bm25}
\end{equation}

Trong đó $k_1$ kiểm soát độ bão hòa tần suất từ, $b$ kiểm soát chuẩn hóa độ dài tài liệu, và $\overline{dl}$ là độ dài tài liệu trung bình.

TechStore triển khai Hybrid Search lấy cảm hứng từ BM25 với cách tính điểm đơn giản hóa phù hợp cấu trúc dữ liệu sản phẩm (chi tiết thiết kế ở Chương~3, §\ref{sec:c3_hybrid_search}).

\subsection{Thách thức ngôn ngữ tiếng Việt trong RAG pipeline}
\label{sec:c2_vietnamese_nlp}

Tiếng Việt có đặc điểm ngôn ngữ đặc thù ảnh hưởng trực tiếp đến chất lượng retrieval trong RAG~\cite{nguyen2018vncorenlp}: (1) viết tắt thương hiệu phổ biến (``ip'' cho iPhone, ``ss'' cho Samsung) làm giảm độ tương đồng cosine khi vector hóa; (2) gõ không dấu (``dien thoai'', ``gia bao nhieu'') giảm chất lượng embedding; (3) pha trộn tiếng Việt và thuật ngữ kỹ thuật tiếng Anh (``laptop gaming pin trâu dưới 20 triệu''). Mô hình multilingual-e5~\cite{wang2022text} và Jina v3~\cite{sturua2024jina} xử lý được phần nào nhờ huấn luyện đa ngôn ngữ, nhưng hiệu quả phụ thuộc vào việc query đã được chuẩn hóa hay chưa~\cite{ma2023query}.

TechStore giải quyết bằng bước chuẩn hóa 71 mẫu regex trước khi vector hóa, hoàn thành dưới 1ms (so với 500--2.000ms nếu dùng LLM rewrite).

\subsection{Query Rewriting trong RAG nâng cao}

\textbf{Query Rewriting} là kỹ thuật mở rộng của RAG: sử dụng LLM để phân tích câu hỏi gốc và viết lại thành dạng tối ưu hơn cho việc tìm kiếm~\cite{ma2023query}. Ví dụ, câu hỏi ``ip17 bnh'' được viết lại thành ``iPhone 17 giá bao nhiêu'' (dạng đầy đủ hơn, giúp tìm kiếm chính xác hơn). Trong TechStore, quá trình rewriting chạy \textbf{song song} với hybridSearch qua \texttt{Promise.all}: LLM rewrite và Hybrid Search diễn ra đồng thời, giảm độ trễ so với tuần tự, sau đó kết quả tốt hơn được ưu tiên sử dụng.

\subsection{Lựa chọn kiến trúc vector store}
\label{sec:c2_vector_store}

TechStore lưu trữ vector trên file JSON cục bộ với tìm kiếm cosine similarity tuyến tính ($O(n)$). Với catalog dưới 10.000 sản phẩm, tìm kiếm hoàn thành trong 30--80ms mà không cần dịch vụ ngoài. Khi cần mở rộng, có thể chuyển sang pgvector hoặc Qdrant bằng cách thay thế lớp \texttt{HybridVectorStore}.

\section{Bảng tổng quan tech stack}
\label{sec:c2_tech_summary}

Bảng~\ref{tab:tech_stack} tổng hợp toàn bộ các công nghệ được sử dụng trong dự án TechStore, phân theo nhóm chức năng.

\begin{table}[H]
  \centering
  \caption{Tổng quan công nghệ sử dụng trong hệ thống TechStore}
  \label{tab:tech_stack}
  \begin{tabular}{|l|l|p{6.5cm}|}
    \hline
    \textbf{Nhóm} & \textbf{Công nghệ} & \textbf{Mục đích} \\
    \hline
    \multirow{6}{*}{Backend} & Node.js 22 LTS & Runtime server, non-blocking I/O \\
    \cline{2-3}
    & Express 4 & Web framework, middleware chain \\
    \cline{2-3}
    & Sequelize 6 + MySQL 8 & ORM, migration, giao dịch \\
    \cline{2-3}
    & JWT + bcrypt & Xác thực, băm mật khẩu \\
    \cline{2-3}
    & Zod & Validation schema tự động \\
    \cline{2-3}
    & google-auth-library & Xác minh Google ID token \\
    \hline
    \multirow{4}{*}{Backend -- AI} & Jina Embeddings v3 & Embedding provider ưu tiên \\
    \cline{2-3}
    & HF multilingual-e5 & Embedding provider dự phòng \\
    \cline{2-3}
    & OpenAI-compatible API & LLM generation \\
    \cline{2-3}
    & Vector DB (JSON) & Lưu trữ embedding sản phẩm \\
    \hline
    \multirow{4}{*}{Backend -- Hạ tầng} & Sharp & Xử lý và chuẩn hóa ảnh \\
    \cline{2-3}
    & node-cron & Tác vụ định kỳ (cleanup) \\
    \cline{2-3}
    & Nodemailer & Gửi email OTP và xác nhận \\
    \cline{2-3}
    & Swagger/OpenAPI 3.0 & Tài liệu API tự động \\
    \hline
    \multirow{10}{*}{Frontend} & React 19 + TypeScript 5 & UI framework, kiểu tĩnh \\
    \cline{2-3}
    & Vite 8 & Build tool, HMR tức thì \\
    \cline{2-3}
    & React Router v7 & Client-side routing, lazy loading \\
    \cline{2-3}
    & shadcn/ui (Radix UI + cva) & Component UI accessible, headless \\
    \cline{2-3}
    & TanStack Query & Server state, cache tự động \\
    \cline{2-3}
    & Zustand (6 stores) & Client state \\
    \cline{2-3}
    & Framer Motion & Animation và hiệu ứng UI \\
    \cline{2-3}
    & i18next + react-i18next & Quốc tế hóa (i18n) \\
    \cline{2-3}
    & exceljs 4.4 + leaflet 1.9.4 & Xuất Excel, bản đồ tương tác \\
    \cline{2-3}
    & dayjs 1.11 & Xử lý ngày giờ (thay moment) \\
    \hline
    \multirow{5}{*}{Kiểm thử / CI} & Jest 29 + Supertest & Unit, integration, API test \\
    \cline{2-3}
    & Stryker Mutator & Mutation testing (ngưỡng 70\%) \\
    \cline{2-3}
    & fast-check & Property-based testing \\
    \cline{2-3}
    & GitHub Actions & CI/CD pipeline tự động \\
    \cline{2-3}
    & Husky + lint-staged & Pre-commit hooks \\
    \hline
    \multirow{2}{*}{Thanh toán} & VNPay (HMAC-SHA512) & Cổng thanh toán nội địa \\
    \cline{2-3}
    & MoMo (HMAC-SHA256) & Ví điện tử \\
    \hline
  \end{tabular}
\end{table}

Bảng~\ref{tab:tech_stack} cho thấy TechStore sử dụng stack công nghệ JavaScript/TypeScript thuần nhất trên cả hai tầng backend và frontend, giúp chia sẻ kiến thức giữa hai phía và giảm chi phí chuyển đổi ngữ cảnh trong quá trình phát triển.

\section{Kiểm thử nâng cao}
\label{sec:c2_quality}

Ngoài các phương pháp kiểm thử truyền thống (unit test, integration test, E2E test) theo mô hình Test Pyramid~\cite{fowler2012test_pyramid}, TechStore áp dụng thêm hai kỹ thuật kiểm thử nâng cao để đảm bảo chất lượng vượt mức coverage đơn thuần.

\subsection{Kiểm thử đột biến (Mutation Testing)}

Kiểm thử đột biến là kỹ thuật đánh giá chất lượng bộ test suite bằng cách tự động chèn các thay đổi nhỏ (mutant) vào source code, ví dụ đổi dấu \texttt{>} thành \texttt{>=} hay xóa một điều kiện, sau đó kiểm tra xem test hiện có có phát hiện được sự thay đổi đó không. Một bộ test tốt sẽ ``giết'' (kill) phần lớn các mutant; mutant còn sống sót (survive) cho thấy test chưa kiểm tra đúng behavior mà chỉ đạt coverage về mặt code execution. Kỹ thuật này bổ sung cho coverage-based testing: một test suite có thể đạt 99\% line coverage mà vẫn bỏ sót các kiểm tra logic quan trọng.

TechStore sử dụng Stryker Mutator~\cite{stryker_docs} với ngưỡng mutation score tối thiểu 70\% cho các module nghiệp vụ quan trọng, đảm bảo bộ test không chỉ thực thi code mà còn xác nhận đúng kết quả nghiệp vụ.

\subsection{Kiểm thử dựa trên thuộc tính (Property-based Testing)}

Kiểm thử dựa trên thuộc tính (property-based testing) là phương pháp kiểm thử trong đó người kiểm thử định nghĩa các bất biến (invariant) mà hệ thống phải thỏa mãn với mọi đầu vào hợp lệ, thay vì viết từng test case thủ công. Framework tự động sinh hàng trăm đầu vào ngẫu nhiên để tìm trường hợp vi phạm. Phương pháp này đặc biệt hiệu quả cho các hàm xử lý dữ liệu nghiệp vụ có không gian đầu vào lớn như tính giá, kiểm tra tồn kho hay phân quyền.

TechStore sử dụng fast-check~\cite{fastcheck_docs} để kiểm tra các bất biến nghiệp vụ quan trọng, bao gồm các ràng buộc về tồn kho, giá sản phẩm và trạng thái đơn hàng, bổ sung cho bộ unit test truyền thống theo hướng example-based.

\bigskip
Chương này đã trình bày nền tảng lý thuyết về kiến trúc Modular Monolith, kỹ thuật RAG Hybrid kết hợp dense và sparse retrieval, thách thức ngôn ngữ tiếng Việt và các kỹ thuật kiểm thử nâng cao. Các kiến thức này làm cơ sở cho thiết kế chi tiết hệ thống TechStore ở Chương~3.
